\documentclass[11pt]{article}
\usepackage[margin=1in]{geometry}
\usepackage{amsmath,amssymb,booktabs}
\usepackage[hidelinks]{hyperref}
\title{A finite charged scalar phase and a reciprocal exterior}
\author{Zimmerman programme: CPB-PI-2026-10-16}
\date{16 October 2026}
\begin{document}
\maketitle
\section{A separate candidate, not inherited closure}
A chosen charged scalar phase can provide a selected interior density,
a computed surface tension, and finite charge, while a reciprocal local
response changes from a cubic-dominated interior to a quadratic exterior.
This addresses the previous prescribed-tail obstruction by changing the
phase physics. Fixed-charge stationary spherical solutions are executed
below. The exterior also changes the inverse-square force coefficient;
there is no empirical, dynamic, microscopic or whole-theory closure.

This is a separately dated action candidate. The original positive-square
MeV scalar potential has $E\ge m|Q_c|$ and a repulsive remainder. It does
not supply an energetically protected lump with $E<m|Q_c|$. The attractive
interaction introduced here cannot inherit that bound, the old neutral
production calibration, or its abundance evidence. Metastable or
gravitationally confined bodies are not excluded by this energy observation.
The global host charge $N$ below is distinct from both $Q_c$ and the
auxiliary flux $Q$.

The charged-droplet mechanism is the established Q-ball construction;
Ref.~\cite{PIqball}, section II, gives its scalar conventions, charge and
radial equations. Its field is written $\phi_0 f_{\rm src}/\sqrt2$;
our $f=|\Psi|$ is dimensionful. Its sextic analysis and interface techniques
are prior work, not a novelty claim for this programme. The following
normalizations and the chosen response coupling are derived explicitly.

\section{Particle action, phase selection, and surface}
With signature $(-+++)$ choose
\begin{equation}
 \mathcal L_\Psi=-g^{\mu\nu}\partial_\mu\Psi^*\partial_\nu\Psi-V(f),
 \quad V=m_h^2 f^2-\lambda f^4+g_6 f^6,
 \quad \lambda,g_6>0,\quad m_h^2>\lambda^2/(4g_6).
 \label{pi:action}
\end{equation}
Here $[g_6]={\rm mass}^{-2}$; a sextic EFT needs ultraviolet control.
Define
\begin{equation}
 f_0^2=\lambda/(2g_6),\qquad
 \omega_0^2=m_h^2-\lambda^2/(4g_6)>0.
\end{equation}
The identity
\begin{equation}
 V=\omega_0^2 f^2+g_6 f^2(f^2-f_0^2)^2
\end{equation}
proves nonnegative potential and a minimum of $V/f^2$ at $f_0$ with
$\omega_0<m_h$. For $\Psi=f(r)e^{i\omega t}$ the positive charge convention
is $n=2\omega f^2$, $N=\int d^3x\,n$. In a homogeneous phase
\begin{equation}
 \rho_h=\omega^2 f^2+V,\qquad p_h=\omega^2f^2-V.
\end{equation}
At coexistence $\omega=\omega_0$, $f=f_0$, $p_h=0$ and the selected
charge density is $n_0=2\omega_0 f_0^2$. This is matter with a temporal
phase current, not vacuum energy with $w=-1$.

The exact planar wall grand potential at coexistence is
$U_0=V-\omega_0^2f^2=g_6 f^2(f^2-f_0^2)^2$.
The first integral $f'^2=U_0$ gives
\begin{equation}
 \frac{f^2(z)}{f_0^2}=\frac{1}{1+e^{2z/\ell}},\qquad
 \ell=\frac{1}{\sqrt{g_6}f_0^2},\qquad
 \sigma=\int_0^{f_0}2\sqrt{U_0}\,df
       =\frac{\sqrt{g_6}f_0^4}{2}.
 \label{pi:wall}
\end{equation}
These are planar coexistence results, not exact finite-radius profiles.
A thin spherical droplet has $p_h\simeq2\sigma/R$ and
$\omega-\omega_0\simeq2\sigma/(n_0R)$; its radius scales as
$(3N/(4\pi n_0))^{1/3}$. Charge fixes size, not the baryonic source by itself.

For the larger homogeneous root, $u=f^2$ obeys
$\omega^2=m_h^2-2\lambda u+3g_6u^2$. With $a=-\lambda+3g_6u>0$,
\begin{equation}
 \frac{dp_h}{d\omega}=n,\qquad
 c_s^2=\frac{n}{\omega\,dn/d\omega}
       =\frac{ua}{ua+\omega^2}.
\end{equation}
This is positive and below unity. At the example's coexistence point
it is $1/4$. The scalar action supplies its own material stress and
current; no thermodynamic identification with the previous label density
$b$ or occupied nodal fermions is assumed.

\section{One reciprocal static phase-response energy}
In a local rest-frame static model choose
\begin{align}
 E_{\rm stat}&=\int d^3x\left[
 |\nabla f|^2+\omega^2 f^2+V(f)
 +\frac C2|\nabla\varphi-\boldsymbol P/C|^2
 +W(f,p)+q_b\rho_b\varphi\right],\\
 W(f,p)&=\frac{K_0}{2}(1-s)^2p^2+\frac\beta3s p^3,
 \quad s=f^2/f_0^2,\quad p=|\boldsymbol P|,
 \quad C,K_0,\beta>0.
 \label{pi:energy}
\end{align}
In four-dimensional mass units $[f]={\rm mass}$, $[p]={\rm mass}^2$,
$[C]=[K_0]=1$, $[\beta]={\rm mass}^{-2}$ and $[q_b]={\rm mass}^{-1}$
if $\rho_b$ is a mass-energy density and $[\varphi]={\rm mass}$.
The switching functions are new choices; their critical zero and coefficients
have not been matched or radiatively protected. This is a static model,
not yet the covariant polarization/metric action.

Let $\boldsymbol D=C\nabla\varphi-\boldsymbol P$ with
$\nabla\cdot\boldsymbol D=q_b\rho_b$. For a spherical source,
$D=q_b M_b(r)/(4\pi r^2)$ and
\begin{equation}
 W_p=D/C\equiv d_r,\quad
 Kp+\beta s p^2=d_r,\quad K=K_0(1-s)^2,\qquad
 p=\frac{2d_r}{K+\sqrt{K^2+4\beta s d_r}}.
 \label{pi:root}
\end{equation}
The root is unique at positive $D$ and extends continuously to $s=0$.
At $s=1$, $p=\sqrt{d_r/\beta}$ is cubic; at $s=0$,
$p=d_r/K_0$ is quadratic. The source-centered radial field is regular
through the finite wall in the executed solutions.

At fixed source the correct reduced energy contribution is
$W-d_rp$, not $W$ alone. Indeed, integrate the source term by parts and
use $\nabla\varphi=(\boldsymbol D+\boldsymbol P)/C$ to obtain
$-D^2/(2C)+W-Dp/C$. The first term is independent of the host and omitted
from the reported reduced host energy. Varying $W-d_rp$ at its stationary
$p$ retains $W_f$ and supplies the reciprocal equation
\begin{equation}
 f''+\frac2r f'=\frac12 V_f-\omega^2f+\frac12 W_f,
 \quad
 W_f=\left[K_0(s-1)p^2+\frac\beta3p^3\right]\frac{2f}{f_0^2}.
 \label{pi:bvp}
\end{equation}
Thus the phase is varied with its response backreaction; the interior
is not imposed as a sharp prescribed density step.

\section{Executed finite-charge spherical branches}
Choose dimensionless units $m_h=\lambda=g_6=1$,
$f_0^2=1/2$, $\omega_0=\sqrt{3/4}$, $\ell=2$, $\sigma=1/8$,
$C=1$, $K_0=0.1$, $\beta=1$. Prescribe only the baryons through
$d_r=S r/(1+r^2)^{3/2}$ with $S=0,0.01,0.05$.
For $R_*=20,40$, fix the charge
$N=(4\pi/3)n_0R_*^3=29020.78983,232166.31862$.
Solve \eqref{pi:bvp} together with $N'(r)=8\pi\omega r^2f^2$,
regular $f'(0)=0$, $N(0)=0$, $f(R_{\rm box})=0$ and the specified
$N(R_{\rm box})$. The unknown $\omega$ is solved, not imposed.
$R_{\rm box}=R_*+50$. The singular radial term is handled explicitly
by the collocation solver. The box boundary approximates the infinite
localized scalar tail; amplitudes five units inside it are below $10^{-10}$.

Two meshes/tolerances and independent 8001-point charge/energy quadrature
were used. Representative refined outputs are
\begin{center}
\begin{tabular}{rrrrr}
\toprule
$R_*$ & $S$ & $\omega$ & $E_{\rm red}/(m_h N)$ & max interior $K/(\beta s p)$\\
\midrule
20 & 0 & 0.881021309 & 0.889219746 & ---\\
20 & 0.01 & 0.881021210 & 0.889219080 & 0.029182\\
20 & 0.05 & 0.881020297 & 0.889212062 & 0.012941\\
40 & 0 & 0.873382644 & 0.877254306 & ---\\
40 & 0.01 & 0.873382632 & 0.877254196 & 0.017948\\
40 & 0.05 & 0.873382510 & 0.877253070 & 0.007986\\
\bottomrule
\end{tabular}
\end{center}
Here $E_{\rm red}=E_h+\int(W-d_rp)$; the omitted source-only energy
is finite. $E_h/(m_hN)<1$ also holds separately, but that comparison
alone is not a proof of full sourced stability. Half-coexistence-density
radii are approximately $19.504,39.533$. The tested interior is
$5<r<0.75R_{1/2}$, not the wall. Its quadratic/cubic force ratio remains
below $0.1$ in all four sourced refined solutions, and the resulting
annulus is finite. The exterior $p\sim r^{-2}$ gives
$4\pi r^3W\sim r^{-1}$; both scalar and response tails have finite energy
in this local candidate. The exterior is an analytic asymptotic matching
after the negligible host tail, not a microscopic nodal-band calculation.

Charge errors are below $4\times10^{-10}$ and mesh changes in
$\omega,E_{\rm red}/N$ below $4\times10^{-9}$. A separate fixed-charge
scaling identity,
\begin{equation}
 2T-3\lambda I_4+6g_6 I_6+
 \int d^3x\,W_f(\tfrac32 f+r f')=0,
 \quad I_j=\int d^3x f^j,\quad T=\int d^3x f'^2,
\end{equation}
checks stationarity without relying only on the solver's residual.
All normalized residuals are below $5\times10^{-9}$.

\section{Finite radial curvature test}
Eliminate $\omega$ at fixed $N$ to obtain $N^2/(4I_2)$ and eliminate
static $p$ at fixed source. The envelope curvature is
\begin{equation}
 H_{ff}=W_{ff}-\frac{W_{fp}^2}{W_{pp}},\qquad W_{pp}=K+2\beta s p.
\end{equation}
For radial perturbations $\delta f=v/r$, the second variation is the
quadratic form of
\begin{equation}
 \mathcal H_0=-2\partial_r^2+V_{ff}-2\omega^2+H_{ff}
 +\frac{32\pi\omega^2}{I_2}|rf\rangle\langle rf|,
 \label{pi:hessian}
\end{equation}
with Dirichlet $v$ at zero and the box boundary and the radial inner
product $\int dr$. The positive rank-one term is required by fixed charge.
The static $S=0$, $p=0$ curvature is understood by the continuous zero
response limit.

Centered differences at 800 and 1600 interior points and three lowest
symmetric eigenpairs per profile yield positive sampled radial amplitude
curvature for all six refined branches. The smallest eigenvalues at 1600
points span $0.398090$--$0.398097$ for $R_*=20$ and
$0.216946$--$0.216949$ for $R_*=40$. Eigenvector residuals are below
$4\times10^{-10}$ and resolution shifts below $1.1\times10^{-4}$.
This is finite static numerical evidence, not a continuum spectral proof,
angular or nonlinear stability result, or a polarization kinetic/causality
certificate. The phase and charge can still have dynamical obligations
outside this amplitude test.

\section{An exterior-force cost and remaining bridge}
The desired finite exterior is not invisible to gravity. Outside,
\begin{equation}
 \varphi'=\frac DC\left(1+\frac{1}{CK_0}\right),\qquad
 G_{\rm scalar,out}=\frac{q_b^2}{4\pi C}
               \left(1+\frac{1}{CK_0}\right).
\end{equation}
The interior bare inverse-square coefficient is $q_b^2/(4\pi C)$;
the example's exterior coefficient is eleven times larger. Einstein
attraction must be added once in the actual measured dictionary. This is
not itself an observational exclusion, but it forbids importing the old
Newtonian calibration or a lensing verdict unchanged.

There is a conditional thin-wall tradeoff. In the selected units the
Laplace-pressure shift gives $s-1\simeq1/R$, hence $K\sim K_0/R^2$.
Near $r\sim R$ on a putative cubic branch,
$K/(\beta s p)\sim K_0/(\sqrt{\beta S}R)$. Demanding small exterior
scalar enhancement $1/(CK_0)<\eta$ and interior ratio below $\epsilon$
therefore requires parametrically
$R\gtrsim1/(C\eta\epsilon\sqrt{\beta S})$. This estimate assumes a
thin weakly sourced phase and is not a universal bound or a data fit.
It identifies a cost of simply raising $K_0$ to hide the exterior.

The new host needs a charge inventory and formation mechanism. A neutral
population could contain charged droplets of opposite signs, but their
production, separation, collisions and annihilation have not been derived.
Sharing $Q$-dependent scales would not by itself create that inventory.
A finite phase of this type is a candidate particle-level host, not an
identification of the previous MeV excitations or the meV nodal sea.
The sextic cutoff, critical matching, covariant polarization dynamics,
complete stress/lensing, parameter origins, flux/history residual, vacuum
magnitude, abundance and cosmology remain open. No complete theory or
global novelty claim follows from these finite stationary solutions.

\section{Evidence}
\texttt{phase\_interface\_probe\_run\_001}: 87 controls including exact
wall identities, twelve BVP solutions, charge/virial/root checks and
refinement comparisons. \texttt{phase\_radial\_curvature\_run\_001}:
twelve finite eigenpair residual checks, with curvature signs retained as
scientific flags. Scripts, input profiles, bounds and raw outputs are pinned
by their manifests. Both current manifests passed freshness validation.
This report is self-reviewed, without an independent referee.

\begin{thebibliography}{9}
\bibitem{PIqball} J. Heeck, A. Rajaraman, R. Riley and C. B. Verhaaren,
\emph{Understanding Q-Balls Beyond the Thin-Wall Limit},
arXiv:2009.08462v2 (9 February 2021),
\url{https://arxiv.org/html/2009.08462v2}, section II, equations (1)--(8).
The established construction is used with the explicit normalization
translation above, not imported as proof of this coupled candidate.
\end{thebibliography}
\end{document}
